\SourceLeaf{236}%
\noindent or we can set
\[
\cot(431)=\alpha,\qquad\cot(234)=\beta
\]
to obtain
\[
\alpha^{2}\cot(341)^{2}-\beta^{2}\left(\frac{\operatorname{tg}i.24}{i}\right)^{2}=1.
\]
Let the tetrahedron $1234$, whose volume we call $\Delta$, now undergo an infinitesimal increase in volume $1241'2'4'=\partial\Delta$ by extending the edge $31$ by the infinitesimal segment $11'=d(13)$ and passing through $1'$, perpendicular to $31'$, the plane $1'2'4'$, which intersects the edges $32$ and $34$ at $2'$ and $4'$, respectively. The angles at vertex $3$, and hence the quantities $\alpha$ and $\beta$, remain unchanged. The angle $(341)$ changes to the angle
\[
(34'1')=(341)+d(341)
\]
and, as is immediately clear from the quadrilateral $11'4'4$ with infinitesimal base $11'=d(13)$ and right angles at $1$ and $1'$, we have
\[
d(341)=-\frac{\sin(i.14)}{i}.d(13).
\]
The increase in volume $1241'2'4'=\partial\Delta$ is bounded laterally by the triangles $124$ and $1'2'4'$, whose planes are both perpendicular to $11'$. Hence (cf. p.~233 of this volume) we have
\[
\partial\Delta=\tfrac{1}{2}d(13).(24).\frac{\sin(i.14)}{i},
\]
consequently
\[
\partial\Delta=-\tfrac{1}{2}(24)d(341),
\]
and, apart from the factor $\tfrac{1}{2}$, this is \textsc{Gauss}'s formula. Thus, if we set
\[
(24)=x,\qquad(341)=\varphi,
\]
the expression for the volume of the tetrahedron $1234$ is
\[
\Delta=-\tfrac{1}{2}\int_{0}^{x}x\frac{d\varphi}{dx}\,dx;
\]
for $\dfrac{d\varphi}{dx}$, the equation
\[
\alpha^{2}\cos\varphi^{2}-\beta^{2}\left(\frac{\operatorname{tg}ix}{i}\right)^{2}=1
\]
is to be used in the calculation.
\begin{flushright}\textsc{Stäckel.}\end{flushright}
\par\endgroup
